\section*{Hoofdstuk \ref{ch:g6:matter-from-food} --- Materie uit voedsel}

\begin{solution}{exo:g6:matter-from-food:1}
\omterm{def:g5:food-webs:roles}{Consumenten} zetten de \omterm{def:g6:plants-are-producers:matter}{organische materie} die ze eten om in hun eigen
\omterm{def:g6:plants-are-producers:matter}{organische materie}. \omterm{def:g5:food-webs:roles}{Producenten} maken nieuwe \omterm{def:g6:plants-are-producers:matter}{organische materie} uit
zuiver minerale grondstoffen, met de \omterm{prop:g3:balanced-diet:jobs}{energie} van het licht.
\end{solution}

\begin{solution}{exo:g6:matter-from-food:2}
De \omterm{def:g4:journey-of-food:digestion}{vertering} haalt het gras uit elkaar tot \omterm{def:g4:journey-of-food:digestion}{voedingsstoffen}
(\cref{ch:g4:journey-of-food}); het \omterm{prop:g4:journey-of-food:absorb}{bloed} bezorgt ze overal
(\cref{ch:g4:heart-and-blood}); de cellen van het lichaam zetten ze weer
in elkaar tot koeienmaterie, en delen zich terwijl het lichaam groeit
(\cref{ch:g5:first-look-at-cells}).
\end{solution}

\begin{solution}{exo:g6:matter-from-food:3}
Groei (de kilo's die het kalf erbij krijgt); vernieuwingen (verharende
\omterm{def:g1:animals-around-us:covering}{vacht}, nieuwe \omterm{def:g1:animals-around-us:covering}{veren}, geweien); herstel (een geheeld \omterm{def:g3:bones-and-muscles:skeleton}{bot}); producten
(\omterm{def:g2:animals-grow-up:born}{melk}, eieren, wol).
\end{solution}

\begin{solution}{exo:g6:matter-from-food:4}
Nooit opgenomen --- eruit als ontlasting; verbrand met de zuurstof uit
de adem voor beweging en warmte --- eruit als uitgeademd gas; en de rest,
gebouwd tot groei en producten.
\end{solution}

\begin{solution}{exo:g6:matter-from-food:5}
Omdat de \qty{25}{kg} \omterm{def:g2:animals-grow-up:born}{melk} van elke dag \omterm{def:g6:plants-are-producers:matter}{organische materie} is die ze uit
het voer van diezelfde dag moet herbouwen; \omterm{def:g2:animals-grow-up:born}{melk} is productie, en productie
is herbouwd voedsel --- geen invoer, geen uitvoer.
\end{solution}

\begin{solution}{exo:g6:matter-from-food:6}
\emph{In} de kip in elkaar gezet, uit de \omterm{def:g4:journey-of-food:digestion}{voedingsstoffen} van haar graan ---
maar de \omterm{def:g6:plants-are-producers:matter}{organische materie} zelf is vóór haar gefabriceerd, in de \omterm{def:g1:plants-around-us:parts}{bladeren}
van de graanplant. De kip zette om; de \omterm{def:g1:plants-around-us:plant}{plant} produceerde.
\end{solution}

\begin{solution}{exo:g6:matter-from-food:7}
Bij elke pijl wordt het meeste van de gegeten materie voor het leven
verbrand of gaat het als ontlasting verloren, en wordt maar een
bescheiden fractie lichaam van de eter --- het enige deel dat de volgende
verdieping kan eten. Elke verdieping heeft daardoor maar een fractie van
de materie eronder: de piramide wordt door rekenkunde smaller.
\end{solution}

\begin{solution}{exo:g6:matter-from-food:8}
In: vlees of brokjes (en af en toe een muis). Gebouwd: haar groei als
kitten, de \omterm{def:g1:animals-around-us:covering}{vacht} die ze bij elke rui vernieuwt, herstel na schrammen.
Uit: poep in de tuin, en het verbrande deel dat elke sprong, elk gespin
en elk warm dutje aandreef.
\end{solution}

\begin{solution}{exo:g6:matter-from-food:9}
Wat verschilt zijn de bron en de uitrusting: gras tegenover konijn,
maalkiezen-en-enorme-darm tegenover grijpers-en-korte-darm. Wat identiek
is, is het vak zelf: allebei halen ze gegeten \omterm{def:g6:plants-are-producers:matter}{organische materie} uit
elkaar tot \omterm{def:g4:journey-of-food:digestion}{voedingsstoffen} en bouwen ze die opnieuw op als hun eigen stof.
\end{solution}

\begin{solution}{exo:g6:matter-from-food:10}
Het meeste vertrok als het verbrande deel --- dat een jaar rennen, denken
en op \qty{37}{\degreeCelsius} blijven aandreef --- en een deel werd
nooit opgenomen en vertrok als ontlasting. Alleen de kleine rest, de
\qty{3}{kg}, werd nieuw lichaam: de boekhouding van het kind klopt net
als die van de koe.
\end{solution}

\begin{solution}{exo:g6:matter-from-food:11}
Het lichaam van de slak haalde de minerale ingrediënten van het huisje
uit haar voedsel en haar \omterm{def:g6:exploring-our-environment:components}{omgeving}, vervoerde ze, en legde ze laag voor
laag in haar eigen bouwwerk neer. De stof is mineraal; het \emph{bouwen}
is het werk van het levende \omterm{def:g1:animals-around-us:animal}{dier} --- een gemaakt ding, van mineraal
materiaal.
\end{solution}

\begin{solution}{exo:g6:matter-from-food:12}
Via het varken wordt het meeste van de materie van de aardappelen aan het
eigen leven van het varken besteed --- verbrand voor warmte en beweging,
verloren als mest --- en komt maar een fractie als varkensvlees terug.
Rechtstreeks gegeten slaan de aardappelen de uitgaven van het varken
over. Het verschil ging de boekhouding van het varken in, niet de tafel op.
\end{solution}

\begin{solution}{pb:g6:matter-from-food:1}
\textbf{1.} De grassen van het weiland en de \omterm{def:g1:plants-around-us:plant}{planten} van de moestuin (en
eventuele fruitbomen). Grondstoffen: water en \omterm{prop:g4:what-plants-need:needs}{minerale zouten} uit de
bodem, koolstofdioxide uit de lucht. \omterm{prop:g3:balanced-diet:jobs}{Energie}: licht.

\textbf{2.} In de eigen \omterm{def:g1:plants-around-us:parts}{bladeren} van de grasplanten --- elke gram:
\omterm{def:g5:food-webs:roles}{producenten} fabriceren ter plaatse, overdag.

\textbf{3.} Geef terug wat je leent: de verterende mest vult de \omterm{prop:g4:what-plants-need:needs}{minerale
zouten} van de tuinbodem aan en vervangt wat de oogsten afvoeren.

\textbf{4.} Omdat de koe een \omterm{def:g5:food-webs:roles}{consument} is: haar vak heeft
\emph{organische} materie nodig om uit elkaar te halen en te herbouwen.
\omterm{prop:g4:what-plants-need:needs}{Minerale zouten} en water zijn grondstoffen van een \omterm{def:g5:food-webs:roles}{producent}; zij heeft
geen machinerie om lichaam uit mineralen te bouwen --- dat is precies het
vak dat haar ontbreekt.

\textbf{5.} Een deel gaat er onopgenomen doorheen --- de koeienvlaaien;
een groot deel wordt met ingeademde zuurstof verbrand om lopen, kauwen en
warm blijven aan te drijven; de rest wordt \omterm{def:g2:animals-grow-up:born}{melk} en koe (en de dagbalans
klopt).

\textbf{6.} In: graan. Gebouwd: eieren, haar eigen onderhoud, en --- bij
de rui --- een volledig nieuw verenkleed. Uit: mest en het verbrande
deel. Tijdens de rui gaat de regel ``gebouwd'' naar de \omterm{def:g1:animals-around-us:covering}{veren}; bij een
gelijkblijvende invoer moet de post eieren dus een tijdlang krimpen.

\textbf{7.} Gras $\rightarrow$ kip (via graan, zelf een plantenproduct)
$\rightarrow$ kuiken $\rightarrow$ buizerd; of aan de weilandkant: gras
$\rightarrow$ koe. Materie werd alleen bij de groene schakels
\emph{gemaakt} (gras, graanplanten); elke dierlijke schakel --- kip,
kuiken, buizerd, koe --- zette haar alleen om.

\textbf{8.} Brood: de \omterm{def:g1:plants-around-us:parts}{bladeren} van de tarweplant. Gekookt \omterm{def:g2:animals-grow-up:born}{ei}: kip uit
graan --- de \omterm{def:g1:plants-around-us:parts}{bladeren} van de graanplant. \omterm{def:g2:animals-grow-up:born}{Melk}: koe uit gras --- de
grasbladeren. Drie sporen, drie \omterm{def:g1:plants-around-us:parts}{bladeren}, allemaal in het licht.

\textbf{9.} Elke verdieping geeft maar een fractie door: van de
\qty{10000}{kg} gras bouwt een konijnenpopulatie maar haar bescheiden
fractie als lichaam; daarvan bouwen vossen weer maar een fractie. Twee
verdiepingen verlies krimpen een weiland vol gras tot een hol vol vossen
--- haar schatting is de rekenkunde van de piramide.

\textbf{10.} Rechtstreeks gegeten bereikt de opbrengst van de tuin het
gezin zonder consumentenverdieping ertussen: niets ervan wordt in het
leven van een \omterm{def:g1:animals-around-us:animal}{dier} verbrand of in zijn mest verloren. De weilandroute
betaalt eerst alle uitgaven van de koe, dus voedt elke vierkante meter
minder mensen.

\textbf{11.} Het gras stopt met produceren, dus valt de invoer van de koe
weg: eerst krimpt de \omterm{def:g2:animals-grow-up:born}{melk} (productie is herbouwd voedsel), daarna
verbrandt ze haar eigen vetreserves en verliest ze gewicht --- de regel
``gebouwd'' wordt negatief. Zonneschijn is \omterm{prop:g3:balanced-diet:jobs}{energie}, geen stof: zonder
water heeft het gras niets om mee te bouwen, dus kan het weiland niet uit
de hemel gered worden.

\textbf{12.} \omterm{def:g5:food-webs:roles}{Producenten}: de \omterm{def:g1:plants-around-us:plant}{planten} van het bedrijf maken al zijn
\omterm{def:g6:plants-are-producers:matter}{organische materie} uit minerale grondstoffen en licht. \omterm{def:g5:food-webs:roles}{Consumenten}: koe,
kippen en gezin zetten die materie om, verbranden het meeste en bouwen de
rest. Ontbreekt: de \omterm{def:g5:food-webs:roles}{afbrekers} --- die de mest, de koeienvlaaien en de
herfstresten weer tot \omterm{def:g6:plants-are-producers:matter}{minerale materie} moeten maken; zij zijn het
onderwerp van het volgende hoofdstuk.
\end{solution}
